\documentclass[10pt,a4paper]{amsart}
\usepackage{amsmath,amssymb,amsthm,graphicx}
% Retained minimal preamble for the Stacks Project chapter sites-cohomology.tex (Cohomology on Sites).
% Verbatim excerpt of the chapter preamble: Xy-pic setup (native preamble lines 14-18) and the native
% theorem declarations with the mathematical macros (native preamble lines 154-217). Copied mechanically;
% nothing added, removed or redefined. The wrapper class and the amsmath/amssymb/amsthm/graphicx base
% packages are supplied by the assembler, so the class/externaldocument part of the native preamble is omitted.
% For commutative diagrams we use Xy-pic
\usepackage[all]{xy}

% We use 2cell for 2-commutative diagrams.
\xyoption{2cell}
% Theorem environments.
%
\theoremstyle{plain}
\newtheorem{theorem}[subsection]{Theorem}
\newtheorem{proposition}[subsection]{Proposition}
\newtheorem{lemma}[subsection]{Lemma}

\theoremstyle{definition}
\newtheorem{definition}[subsection]{Definition}
\newtheorem{example}[subsection]{Example}
\newtheorem{exercise}[subsection]{Exercise}
\newtheorem{situation}[subsection]{Situation}

\theoremstyle{remark}
\newtheorem{remark}[subsection]{Remark}
\newtheorem{remarks}[subsection]{Remarks}

\numberwithin{equation}{subsection}

% Macros
%
\def\lim{\mathop{\mathrm{lim}}\nolimits}
\def\colim{\mathop{\mathrm{colim}}\nolimits}
\def\Spec{\mathop{\mathrm{Spec}}}
\def\Hom{\mathop{\mathrm{Hom}}\nolimits}
\def\Ext{\mathop{\mathrm{Ext}}\nolimits}
\def\SheafHom{\mathop{\mathcal{H}\!\mathit{om}}\nolimits}
\def\SheafExt{\mathop{\mathcal{E}\!\mathit{xt}}\nolimits}
\def\Sch{\mathit{Sch}}
\def\Mor{\mathop{\mathrm{Mor}}\nolimits}
\def\Ob{\mathop{\mathrm{Ob}}\nolimits}
\def\Sh{\mathop{\mathit{Sh}}\nolimits}
\def\NL{\mathop{N\!L}\nolimits}
\def\CH{\mathop{\mathrm{CH}}\nolimits}
\def\proetale{{pro\text{-}\acute{e}tale}}
\def\etale{{\acute{e}tale}}
\def\QCoh{\mathit{QCoh}}
\def\Ker{\mathop{\mathrm{Ker}}}
\def\Im{\mathop{\mathrm{Im}}}
\def\Coker{\mathop{\mathrm{Coker}}}
\def\Coim{\mathop{\mathrm{Coim}}}

% Boxtimes
%
\DeclareMathSymbol{\boxtimes}{\mathbin}{AMSa}{"02}

%
% Macros for moduli stacks/spaces
%
\def\QCohstack{\mathcal{QC}\!\mathit{oh}}
\def\Cohstack{\mathcal{C}\!\mathit{oh}}
\def\Spacesstack{\mathcal{S}\!\mathit{paces}}
\def\Quotfunctor{\mathrm{Quot}}
\def\Hilbfunctor{\mathrm{Hilb}}
\def\Curvesstack{\mathcal{C}\!\mathit{urves}}
\def\Polarizedstack{\mathcal{P}\!\mathit{olarized}}
\def\Complexesstack{\mathcal{C}\!\mathit{omplexes}}
% \Pic is the operator that assigns to X its picard group, usage \Pic(X)
% \Picardstack_{X/B} denotes the Picard stack of X over B
% \Picardfunctor_{X/B} denotes the Picard functor of X over B
\def\Pic{\mathop{\mathrm{Pic}}\nolimits}
\def\Picardstack{\mathcal{P}\!\mathit{ic}}
\def\Picardfunctor{\mathrm{Pic}}
\def\Deformationcategory{\mathcal{D}\!\mathit{ef}}
\usepackage{cleveref}
\crefname{originallemma}{original lemma}{original lemmas}
\Crefname{originallemma}{Original lemma}{Original lemmas}
\makeatletter
\newlabel{lemma-pullback-K-flat}{{06YC}{source}{}{sourcepointer-1}{}}
\newlabel{lemma-pullback-K-flat@cref}{{[originallemma][0][]06YC}{[1][0][]source}{}{}{}}
\makeatother
\crefname{originalsection}{original section}{original sections}
\Crefname{originalsection}{Original section}{Original sections}
\makeatletter
\newlabel{section-compare}{{073L}{source}{}{sourcepointer-2}{}}
\newlabel{section-compare@cref}{{[originalsection][0][]073L}{[1][0][]source}{}{}{}}
\makeatother
\title{The cartesianity condition for a comparison with the chaotic topology makes epsilon\textasciicircum{}* and R epsilon\_* inverse equivalences onto the quasi-coherent objects}
\author{The Stacks Project Authors}
\date{stacks-project:sites-cohomology.tex:local-snapshot-20261011; GNU FDL 1.2 or later}
\begin{document}
\maketitle
\noindent\textit{Editorial note.} All delivered mathematical text is the author's own, copied from the pinned Stacks Project chapter file sites-cohomology.tex; no statement or proof step was rewritten, reordered or removed.\par\medskip
\noindent\textit{Editorial note.} The counted proof is the author's complete argument: the proof is the author's complete argument: the fullness reduction, the choice of a K-flat representative, the Leray computation of the round trip over an arbitrary object U, the identification of the resulting tensor-up complex using quasi-coherence of M, and the appeal to condition (2) are all written out.\par\medskip
\noindent\textit{Editorial note.} The label lemma-cartesion-plus-topology names the counted statement, and the proof environment immediately following it in the source is delivered verbatim as its proof.\par\medskip
\noindent\textit{Editorial note.} The statement and proof are copied from the pinned chapter file, including the footnote explaining the quasi-coherence condition; the two cited inputs are kept as external pointers with their canonical Stacks tags.\par\medskip
\noindent\textit{Editorial note.} The preamble is the chapter's own retained minimal preamble (sources/stacks-sites-cohomology/retained-preamble.tex), a verbatim excerpt of the native preamble: the Xy-pic setup and the native theorem declarations with the mathematical macros.\par\medskip
\noindent\textit{Allowed standard prerequisites (editorial).}\begin{itemize}
\item The comparison situation with a coarser topology tau and the functors epsilon\textasciicircum{}*, R epsilon\_*. Same chapter, the section comparing the topology tau with tau' and the observations made there about epsilon\textasciicircum{}* and R epsilon\_*. Supplies the functors whose equivalence is asserted and the observations used without further mention.
\item Quasi-coherent objects QC(O) of D(C\_tau', O). Same chapter, the definition of QC(O) by the cartesianity of the derived global sections functor. Supplies the subcategory appearing on both sides of the asserted equivalence and the condition used at the end of the proof.
\item K-flat resolutions and their pullbacks. Same chapter, the lemma stating that the pullback of a K-flat complex with flat terms is again K-flat, used to see that M\textasciicircum{}bullet(U) is K-flat.. Supplies the choice of the K-flat representative M\textasciicircum{}bullet used in the computation.
\end{itemize}\medskip
\ifdefined\hypertarget\hypertarget{sourcepointer-1}{}\fi
\noindent\textit{Source reference.} Original lemma 06YC: sources/stacks-sites-cohomology/pinned-original-source.tex, source lines 3508--3508. Prerequisite pointer invoked at proof line(s) 11661; it names a standard fact of the same chapter that the counted argument uses, and every step of the counted proof built on it is retained verbatim in the delivered body.\par\medskip
\ifdefined\hypertarget\hypertarget{sourcepointer-2}{}\fi
\noindent\textit{Source reference.} Original section 073L: sources/stacks-sites-cohomology/pinned-original-source.tex, source lines 6177--6177. Prerequisite pointer invoked at proof line(s) 11611; it names a standard fact of the same chapter that the counted argument uses, and every step of the counted proof built on it is retained verbatim in the delivered body.\par\medskip
\begin{lemma}
\label{lemma-cartesion-plus-topology}
Let $\epsilon : (\mathcal{C}_\tau, \mathcal{O}_\tau) \to
(\mathcal{C}_{\tau'}, \mathcal{O}_{\tau'})$ be as in
Section \ref{section-compare}. Assume
\begin{enumerate}
\item $\tau'$ is the chaotic topology on the category $\mathcal{C}$,
\item for all $U \in \Ob(\mathcal{C})$ and all K-flat complexes of
$\mathcal{O}(U)$-modules $M^\bullet$ the map
$$
M^\bullet \longrightarrow
R\Gamma((\mathcal{C}/U)_\tau,
(M^\bullet \otimes_{\mathcal{O}(U)} \mathcal{O}_U)^\#)
$$
is a quasi-isomorphism (see proof for an explanation).
\end{enumerate}
Then $\epsilon^*$ and $R\epsilon_*$ define mutually quasi-inverse
equivalences between $\mathit{QC}(\mathcal{O})$ and
the full subcategory of $D(\mathcal{C}_\tau, \mathcal{O}_\tau)$
consisting of objects $K$ such that $R\epsilon_*K$ is in
$\mathit{QC}(\mathcal{O})$\footnote{This means that
$R\Gamma(V, K) \otimes_{\mathcal{O}(V)}^\mathbf{L} \mathcal{O}(U)
\to R\Gamma(U, K)$ is an isomorphism for all $U \to V$ in $\mathcal{C}$.}.
\end{lemma}

\begin{proof}
We will use the observations made in Section \ref{section-compare} without
further mention. Since $R\epsilon_*$
is fully faithful and $\epsilon^* \circ R\epsilon_* = \text{id}$,
to prove the lemma it suffices to show that for $M$ in
$\mathit{QC}(\mathcal{O})$ we have $R\epsilon_*(\epsilon^*M) = M$.
Condition (2) is exactly the condition needed to see this.
Namely, we choose a K-flat complex $\mathcal{M}^\bullet$
of $\mathcal{O}$-modules with flat terms representing $M$.
Then we see that $\epsilon^*M$ is represented by the $\tau$-sheafification
$(\mathcal{M}^\bullet)^\#$ of $\mathcal{M}^\bullet$.
Let $U \in \Ob(\mathcal{C})$. By Leray we get
$$
R\Gamma(U, R\epsilon_*(\epsilon^*M)) =
R\Gamma((\mathcal{C}/U)_\tau, (\mathcal{M}^\bullet)^\#|_{\mathcal{C}/U}) =
R\Gamma((\mathcal{C}/U)_\tau, (\mathcal{M}^\bullet|_{\mathcal{C}/U})^\#)
$$
The last equality since
sheafification commutes with restriction to $\mathcal{C}/U$.
As usual, denote $\mathcal{O}_U$ the restriction of $\mathcal{O}$ to
$\mathcal{C}/U$. Consider the map
$$
\mathcal{M}^\bullet(U) \otimes_{\mathcal{O}(U)} \mathcal{O}_U
\longrightarrow
\mathcal{M}^\bullet|_{\mathcal{C}/U}
$$
of complexes of $\mathcal{O}_U$-modules (in $\tau'$-topology).
By our choice of $\mathcal{M}^\bullet$ the complex $\mathcal{M}^\bullet(U)$
is a K-flat complex of $\mathcal{O}(U)$-modules; see
Lemma \ref{lemma-pullback-K-flat} and use that the inclusion
of $U$ into $\mathcal{C}$ defines a morphism of ringed topoi
$(\Sh(pt), \mathcal{O}(U)) \to (\Sh(\mathcal{C}_{\tau'}), \mathcal{O})$.
Since $M$ is in $\mathit{QC}(\mathcal{O})$ we conclude that
the displayed arrow is a quasi-isomorphism. Since sheafification
is exact, we see that the same remains true after sheafification.
Hence
$$
R\Gamma(U, R\epsilon_*(\epsilon^*M)) =
R\Gamma((\mathcal{C}/U)_\tau,
(M^\bullet \otimes_{\mathcal{O}(U)} \mathcal{O}_U)^\#)
$$
and assumption (2) tells us this is equal to
$R\Gamma(U, M) = \mathcal{M}^\bullet(U)$ as desired.
\end{proof}

\end{document}
